\section{Momenta: producción en serie, valor para el cliente y el «bagre»}

La sección sobre Waymo explica cómo se descompone un sistema; la sección sobre Momenta explica cómo los clientes ponen a prueba un sistema. Para Gao Jiyang (高继扬), pasar de Waymo a Momenta no fue saltar de una buena empresa a otra buena empresa, sino pasar del campo de entrenamiento de ingenieros al terreno de la entrega en producción en serie.

Si Waymo es el paraíso de los ingenieros, Momenta es el campo de batalla de la producción en serie. Gao Jiyang eligió Momenta porque quería pasar de los sistemas de conducción autónoma al producto, a los clientes y a la gestión de una empresa; en particular, quería entender cómo un sistema de nivel demo se convierte en un producto que pueda entregarse a un fabricante de automóviles.

\begin{figure}[H]
\centering
\includegraphics[width=0.96\textwidth]{waymo-momenta-contrast.png}
\caption{Waymo y Momenta: dos formas extremas de formación organizacional.}
\end{figure}

\begin{knowledgebox}{Lectura de la figura: qué forma cada organización}
A la izquierda, Waymo forma en ingeniería de sistemas, infra, calidad del código y ejecución de la estrategia a largo plazo; a la derecha, Momenta forma en entrega, comunicación con el cliente, ajuste organizacional y el ciclo de valor comercial. Una enseña cómo funciona un sistema; la otra enseña cómo los clientes usan un sistema, pagan por él y obligan a iterarlo.
\end{knowledgebox}

\subsection{Por qué la producción en serie es una ruta de datos}

La lógica central de Momenta es esta: la conducción autónoma necesita, en última instancia, una gran cantidad de datos reales. Si se depende solo de una flotilla propia, la cobertura de ciudades y escenarios avanza demasiado despacio; si capacidades como la asistencia a la conducción o el estacionamiento se instalan en vehículos de producción en serie, primero se crea valor para el cliente y después se obtienen datos a través de esos vehículos, con lo que se forma un ciclo de datos impulsado por el negocio. Lo decisivo aquí no es «venderle a los fabricantes de automóviles» en sí, sino que la obtención de datos deja de ser un costo puro y pasa a pagarse con el valor que recibe el cliente.

\begin{importantbox}{Ciclo de datos y valor comercial}
La producción en serie no es solo la entrega de un proyecto: hace que el producto llegue a vehículos reales, clientes reales y caminos reales, y así convierte la recolección de datos en una actividad impulsada por el negocio. Cuanto más claro es el valor para el cliente, más sostenible es el retorno de datos.
\end{importantbox}

\subsection{El papel de bagre}

Lo anterior trató la transformación organizacional de Momenta; aquí el enfoque vuelve al propio Gao Jiyang. Lo que se llama bagre no es una etiqueta de personalidad, sino una forma de intervenir en sistemas que se forjó una y otra vez en un entorno de entrega bajo alta presión.

El «bagre» del título de la entrevista se refiere a que en Momenta a Gao Jiyang lo asignaban con flexibilidad a distintos módulos: percepción, localización, estacionamiento, infra, planificación y control, y la producción en serie de NOA. Su propio resumen es: entrar rápido en áreas desconocidas, entender el sistema con una metodología fija, descomponer las tareas, asignar a las personas adecuadas, monitorear la retroalimentación, ampliar cuando la retroalimentación es buena y reducir y ajustar cuando no lo es.

\begin{knowledgebox}{La metodología del bagre}
El bagre no se dedica a agitarlo todo, sino a que el sistema reaccione mejor. Sus acciones incluyen: entrar en un área desconocida, trazar rápidamente un mapa, descomponer tareas, asignar personas, fijar indicadores, observar la retroalimentación y ajustar la organización.
\end{knowledgebox}

\subsection{Resumen de la sección}

Lo que Momenta le dio a Gao Jiyang fue formación en valor para el cliente y en presión organizacional. Le mostró que, para que un equipo técnico se convierta en una empresa de producto, tiene que pasar por la prueba de la entrega, la eliminación, el ajuste y la producción en serie.
